El sistema cumple el objetivo: responde con citas trazables, casi no hace afirmaciones sin respaldo
(98,5\,\% respaldadas) y reconoce cuándo el corpus no alcanza. Corre en una computadora sin GPU, salvo
los modelos de lenguaje, que se usan por API, con un costo de unos 4 centavos de dólar por consulta.

Lo que más influyó en la calidad no fueron los modelos sino la recuperación. La reescritura de la
pregunta en inglés resolvió la diferencia de idioma entre las preguntas y el corpus, y el límite de
fragmentos por artículo pesó más que la cantidad de fragmentos. El límite actual es que en alrededor
de una de cada tres preguntas el contexto no incluye ningún fragmento que responda, y es la causa de
casi todas las respuestas incompletas.

Dos resultados contradijeron lo esperado y se dejaron documentados. El reranker empeoró el orden de
los fragmentos frente a la fusión RRF y no mejoró la cobertura del contexto, así que con este modelo
conviene desactivarlo (está en \texttt{config.json}) o reemplazarlo. El verificador no mejoró la
corrección y su valor es de alarma más que de corrección.

Posibles mejoras:
\begin{itemize}
\item Recuperación en dos niveles: elegir primero los artículos y después buscar dentro de cada uno,
ya que el artículo correcto aparece entre los tres primeros en todas las preguntas.
\item Embeddings multilingües o de dominio biomédico, y un reranker entrenado en búsquedas sobre
PubMed (por ejemplo, MedCPT) en lugar de búsquedas web.
\item Verificación adaptativa: llamar al verificador solo cuando la respuesta tiene muchas citas o
cuando la recuperación tiene puntajes bajos, para bajar costo y latencia.
\item Revisión humana del conjunto de preguntas (pendiente) y más preguntas, para que las diferencias
entre variantes tengan validez estadística. Con 40 preguntas, una diferencia de 0,02 en la
corrección media no es significativa.
\item Ampliar el corpus con IEEE Xplore y con los artículos que hoy solo tienen resumen, y mantenerlo
actualizado con alertas de PubMed.
\item Un conjunto de ataques independiente, porque los ataques y la regla léxica los escribió la
misma persona.
\end{itemize}
